\documentclass[12pt,a4paper]{article}

% ── Packages ──────────────────────────────────────────────────────────────────
\usepackage[margin=2.5cm]{geometry}
\usepackage[numbers,sort&compress]{natbib}
\usepackage{hyperref}
\usepackage{setspace}
\usepackage{parskip}
\usepackage{titlesec}
\usepackage{booktabs}
\usepackage{microtype}
\usepackage{amsmath}
\usepackage{eurosym}
\usepackage{xurl}

% ── Formatting ─────────────────────────────────────────────────────────────────
\onehalfspacing
\hypersetup{colorlinks=true, linkcolor=blue, citecolor=blue, urlcolor=blue}

\titleformat{\section}{\large\bfseries}{\thesection.}{0.5em}{}
\titleformat{\subsection}{\normalsize\bfseries}{\thesubsection}{0.5em}{}

% ══════════════════════════════════════════════════════════════════════════════
\begin{document}

    \begin{titlepage}
        \centering
        \vspace*{2cm}
        {\Large\bfseries RESEARCH PROPOSAL \par}
        \vspace{1.5cm}
        {\large\bfseries AN AUTONOMOUS GEOFENCE-TRIGGERED IMPLANTABLE
            SEDATIVE-DELIVERY DEVICE FOR IMMEDIATE IMMOBILIZATION OF HABITUAL
            WILDLIFE TRESPASSERS \par}
        \vspace{2cm}
        \vfill
    \end{titlepage}

    % ══════════════════════════════════════════════════════════════════════════════
    \begin{abstract}
        \noindent
        Human-wildlife conflict (HWC) represents one of the most urgent and complex challenges
        at the interface of conservation and human livelihoods globally. In regions adjacent to
        game parks and wildlife reserves, conflict animals routinely trespass into human
        settlements, resulting in crop destruction, livestock predation, serious injury, and
        fatalities on both sides of the divide. In India alone, over 500 people are killed annually
        through negative human-elephant interactions, with over 50,000 farming families
        experiencing crop losses amounting to approximately one million hectares per year
        \cite{muliya2026}. In sub-Saharan Africa, nine-year KWS records from Kenya show that
        elephants alone account for 79\,\% of reported HWC incidents, with peak conflict
        occurring during dry-season months as animals disperse from protected areas in search of
        food and water \cite{manoa2020}. Retaliatory killings further endanger conflict-prone
        megafauna such as elephants, big cats, hippopotami, and members of the Big Five.
        Current response frameworks are predominantly reactive, triggered only after an animal
        has already entered human-occupied spaces and caused harm.

        Despite remarkable advances in GPS satellite telemetry-based early warning systems
        \cite{muliya2026}, GPS-enabled virtual geofencing for animal management \cite{gilbert2026},
        low-cost Arduino-based GPS tracking devices \cite{imamoglu2022}, biocompatible implantable
        MEMS drug-delivery micropumps \cite{gutierrez2014}, fully implantable robotic drug
        delivery systems \cite{iacovacci2021}, and autonomous closed-loop drug delivery devices
        \cite{ciatti2024}, no study has yet combined geofence boundary-detection with autonomous
        implantable sedative delivery for wildlife conflict management. This gap leaves
        immobilization response time entirely dependent on human logistics.

        This research aims to design, develop, and validate an autonomous geofence-triggered
        implantable sedative-delivery device that can immediately immobilize habitual wildlife
        trespassers upon crossing a pre-defined virtual boundary, without requiring human
        intervention at the moment of immobilization. The device is intended for subcutaneous
        implantation in recidivist conflict animals, particularly big cats, hippopotami, and members
        of the Big Five. The research will proceed in three phases: device engineering, large-animal
        biocompatibility assessment, and controlled field trials. The ultimate outcome is a
        technology that reduces human casualty rates, minimises crop and property damage,
        decreases retaliatory animal killings, and represents a paradigm shift from reactive to
        proactive HWC management.
    \end{abstract}

    \newpage
    \tableofcontents
    \newpage

    % ══════════════════════════════════════════════════════════════════════════════
    \section{Introduction}

    \subsection{Background}

    Human-wildlife conflict is a globally recognised conservation crisis arising wherever the
    spatial and resource requirements of humans and wildlife overlap. As habitat loss,
    agricultural expansion, and human population growth drive communities alongside
    previously wild landscapes, encounters between people and large, dangerous animals have
    escalated dramatically \cite{muliya2026}. The consequences are profound: local communities
    lose crops, livestock, and lives; wildlife populations face retaliatory persecution that
    threatens conservation goals; and public tolerance for wildlife erodes, undermining the
    social licence for conservation.

    The problem is especially acute around game parks and forest corridors in sub-Saharan
    Africa and South Asia. In Kenya, the severity and economic cost of HWC is well
    documented. A nine-year analysis of Kenya Wildlife Service (KWS) records across Kajiado
    and Laikipia Counties identified 953 valid HWC cases between 2010 and 2018, with an
    overall rising trend peaking in 2017 \cite{manoa2020}. Wildlife threats to human life,
    crop destruction, and livestock predation contributed 65.7\,\%, 21.7\,\%, and 7.7\,\% of
    incidents respectively, with elephants responsible for 79\,\% of all reported cases. Across
    the preceding decade, Kenya recorded 305 human deaths attributable to HWC, while
    between 2014 and 2017 alone, the national toll reached 452 human deaths, 4,555 injuries,
    5,073 crop damage incidents, and 3,012 livestock predation events, representing a combined
    compensation liability of KES\,3.4\,billion (approximately USD\,317\,million) \cite{manoa2020}.
    In Kajiado County specifically, the most common conflict species include elephants,
    hyenas, buffaloes, leopards, baboons, monkeys, snakes, and crocodiles, with elephants
    responsible for 46.2\,\% of conflicts, predominantly through crop raiding \cite{mbiki2025}.

    The socio-economic and ecological consequences of HWC extend beyond the immediate
    incident. Crop destruction, livestock depredation, and infrastructure damage place heavy
    financial burdens on affected communities, exacerbating poverty and discouraging
    conservation participation \cite{mbiki2025}. Retaliatory killings of predators cascade
    ecologically, producing herbivore overpopulation, habitat degradation, and biodiversity
    loss. HWC also generates fear, insecurity, and inter-community tensions that destabilise
    socio-economic landscapes and erode the psychological wellbeing of affected families
    \cite{mbiki2025}. Research further shows that negative perceptions of wildlife : where
    animals are increasingly seen as threats rather than assets : directly undermine conservation outcomes by reducing community support for protected areas and wildlife
    management programmes \cite{basak2023}.

    Understanding the human dimension of HWC is recognised as essential for effective
    management. A systematic review of 124 peer-reviewed studies on urban HWC found that
    managing conflicts requires understanding local attitudes and behaviours, since community
    perceptions of wildlife directly determine whether management interventions will be
    accepted or resisted \cite{basak2023}. Property damage was identified as the most commonly
    reported conflict (25\,\% of studies), followed by danger to human life (12.7\,\%) and
    danger to pets (11\,\%). Critically, Basak et al.\ identify that delayed or insufficient
    compensation actively worsens community attitudes toward wildlife, feeding a cycle of
    resentment and retaliatory killing \cite{basak2023}. This finding is corroborated at the
    policy level in Kenya: the government has acknowledged that between 2013 and 2022,
    KWS paid only KES\,2.7\,billion to HWC victims over nearly a decade, with victims
    sometimes waiting up to eight years for compensation, fuelling hostility toward wildlife
    conservation \cite{kenya2025}. In response, the Kenyan government launched a digitised
    compensation platform in 2025 targeting payment within four months of an incident and
    mandating that all reports of loss be completed within 90\,days \cite{kenya2025}. Supporting
    operational reforms include the recruitment of 1,500 new game rangers, procurement of
    100 vehicles for KWS, and an extension of wildlife fencing by 200\,km \cite{kenya2025}.
    These policy developments reflect the urgency with which governments are now approaching
    HWC, and highlight the need for technological solutions that can complement, rather than
    merely supplement, institutional responses.

    In India, which hosts over 60\,\% of the global wild elephant population, HEC causes over
    500 human deaths and crop losses of approximately one million hectares annually, with
    ex-gratia compensation averaging USD\,4.79\,million per year between 2009 and 2020
    \cite{muliya2026}. In Africa, similar dynamics characterise conflicts involving lions,
    leopards, rhinoceroses, buffaloes, hippopotami, and elephants. The animals most frequently
    involved in recurring conflicts are ``habitual trespassers''-individuals that have learned to
    exploit human-dominated landscapes and repeatedly return despite deterrence efforts.

    A convergence of technologies from disparate fields now offers a compelling pathway to a
    fundamentally new solution. Virtual fencing technology uses GPS-enabled collars to enforce
    digital spatial boundaries on free-ranging animals via conditioned stimuli \cite{gilbert2026}.
    GPS satellite telemetry-based early warning systems have demonstrated that real-time
    tracking of conflict-prone animals can dramatically reduce HWC severity \cite{muliya2026}.
    Low-cost Arduino-based GPS tracking platforms have further demonstrated that
    positionally accurate animal tracking-within 23\,m of industrial-grade GPS benchmarks-
    can be achieved at a fraction of the cost of commercial devices, making real-time wildlife
    monitoring accessible to under-resourced wildlife agencies \cite{imamoglu2022}. Implantable
    MEMS drug delivery devices have established the biocompatibility and reliability of
    miniaturised programmable drug pumps in large animal models \cite{gutierrez2014}. Fully
    implantable robotic systems have demonstrated non-invasive drug reservoir refilling through
    ingestible magnetic capsules, extending device operational life without surgical re-entry
    \cite{iacovacci2021}. Autonomous closed-loop implantable devices have proven that
    sensor-triggered autonomous drug delivery can rescue freely moving large animals from
    life-threatening events within minutes \cite{ciatti2024}. A comprehensive review of the
    BioMEMS field documents the rich landscape of biocompatible materials, miniaturised
    actuation systems, wireless power transfer, and closed-loop sensing architectures that
    underpin next-generation implantable devices \cite{abhinav2025}. Collectively, these bodies
    of work provide strong technical and conceptual foundations for a novel approach to the
    delayed-tranquilization problem.

    \subsection{Problem Statement}

    When encroachment by a conflict animal occurs near game parks and wildlife reserves, it
    currently takes a very long time to safely capture the animal. During this period, serious
    injuries or fatalities to humans, a large number of animal deaths, and extensive crop
    damage are recorded. Wildlife managers must first track the animal using GPS collars and
    implants and then sedate it using dart projectors, blowpipe CO\textsubscript{2}-powered
    rifles, or projectile syringes. This process can take anywhere from a few hours to several
    months when managing multiple conflict animals. Seasonality compounds the challenge:
    HWC incidents peak during dry-season months of July, January, and October, when wildlife
    disperses most aggressively into human settlements, placing maximum pressure on response
    teams that are already stretched across large rangeland areas \cite{manoa2020}. The
    response is necessarily reactive and entirely dependent on human logistics, while casualties
    accumulate. No autonomous device currently exists that can immediately immobilize a
    conflict animal the moment it crosses a designated boundary without requiring the physical
    presence of a trained human responder.

    \subsection{Research Objectives}

    This research addresses the identified gap through the following objectives:

    \begin{enumerate}
        \item To design and develop a GPS geofence-triggered implantable sedative delivery
        device suitable for subcutaneous implantation in habitual wildlife trespassers,
        including big cats, hippopotami, and members of the Big Five.
        \item To evaluate the biocompatibility, structural integrity, and safety profile of
        the device in appropriate large animal models over up to twelve months.
        \item To assess the efficacy of the device in triggering autonomous immobilization
        within an operationally acceptable timeframe upon geofence boundary crossing.
        \item To compare the response time, effectiveness, and safety outcomes of the proposed
        device against conventional dart-based tranquilization methods currently used in
        HWC management.
        \item To identify the ethical, operational, and regulatory considerations governing
        deployment of autonomous implantable immobilization devices in wildlife management.
    \end{enumerate}

    % ══════════════════════════════════════════════════════════════════════════════
    \section{Literature Review}

    \subsection{Human-Wildlife Conflict: Scale, Drivers, and Impacts}

    Human-wildlife conflict is an escalating global conservation and humanitarian challenge.
    At its core, HWC emerges wherever the spatial, temporal, or resource requirements of
    wildlife and humans overlap \cite{muliya2026}. Key drivers include accelerating habitat
    fragmentation, agricultural expansion into wildlife corridors, human population growth in
    peri-protected area zones, and the recovery or expansion of large wildlife populations in
    human-dominated landscapes. The consequences are asymmetrically borne by rural
    communities, who experience crop loss, livestock depredation, physical injury, death, and
    psychological distress.

    In Kenya, the typology and economic magnitude of HWC have been quantified through
    systematic analysis of KWS archival records across Kajiado and Laikipia Counties
    \cite{manoa2020}. Over the 20102018 period, wildlife threats to human life accounted for
    65.7\,\% of 953 documented incidents, with crop destruction contributing 21.7\,\% and
    livestock predation 7.7\,\%. Elephants were responsible for the largest share of all
    incidents (79\,\%, $n=753$), attributable to their large body size, extensive home ranges
    (estimated at 5,2007,790\,km\textsuperscript{2} in the Amboseli landscape), and high
    daily resource requirements of approximately 150\,kg of vegetation and 190\,litres of
    water \cite{manoa2020}. Critically, approximately 80\,\% of elephant range in Kenya lies
    outside protected areas, ensuring frequent contact with human settlements. Baboons were
    the second most conflict-prone species (6.9\,\%, $n=66$), exploiting a wide habitat range
    and omnivorous diet to raid crops, attack livestock, and destroy property \cite{manoa2020}.
    Human attacks in Kajiado County alone resulted in six deaths and 21 injuries over the
    nine-year period, generating compensation claims equivalent to KES\,82.5\,million; the
    responsible species included elephants (43\,\% of combined fatalities and injuries), snakes
    (17\,\%), and buffaloes (17\,\%).

    In South Asia, the Asian elephant (\textit{Elephas maximus}) exemplifies the scale and
    complexity of HWC. India hosts an estimated 25,00030,000 individuals across
    163,000\,km\textsuperscript{2} and is home to the single largest population of wild
    elephants globally, classified as Endangered by the IUCN \cite{muliya2026}. Between 2019
    and 2021, the state of Karnataka alone recorded approximately 35,000 crop-raiding incidents
    and over 50 human fatalities attributable to elephants, with ex-gratia payments totalling
    approximately USD\,4.1\,million to affected communities \cite{muliya2026}. Nationally, over
    50,000 farming families are affected annually, with crop losses estimated at one million
    hectares per year and over 500 human deaths annually from negative human-elephant
    interactions \cite{muliya2026}. More than 100 elephants are killed each year in retaliatory
    actions, creating a feedback loop that drives further illegal persecution \cite{muliya2026}.

    The human dimension of HWC is not limited to physical and economic losses. Research
    drawing on a systematic review of 124 global perception studies identifies that community
    attitudes toward wildlife are a critical determinant of conservation outcomes \cite{basak2023}.
    In urbanising areas adjacent to protected areas, human perceptions of wildlife range from
    appreciation to fear and even hatred, and these perceptions are shaped not only by direct
    experience but by cultural norms, historical encounters, and the perceived fairness of
    government responses to losses \cite{basak2023}. Negative perceptions - where animals are
    seen primarily as threats - undermine community willingness to participate in conservation
    and increase tolerance for retaliatory killings. In Kajiado County, historical patterns of
    Maasai coexistence with wildlife have been disrupted by increasing livestock predation and
    inadequate government response, driving retaliatory killing of predators and a deterioration
    of human-wildlife relations \cite{mbiki2025}. The frequency with which wildlife is perceived
    in urban and peri-urban landscapes, and the speed and fairness with which conflict-related
    losses are compensated, have been identified as key moderators of community tolerance
    \cite{basak2023,mbiki2025}. These findings establish that any technological intervention in
    HWC management must be complemented by systems that address community perception
    and restore confidence in wildlife governance.

    \subsection{Current Mitigation Strategies and Their Limitations}

    Wildlife managers globally have implemented a diverse portfolio of HWC mitigation
    measures. These can be broadly categorised as: physical and chemical barriers such as
    electrified fencing, chili-based deterrents, and beehive fences; compensation and
    ex-gratia schemes; community-based vigilance and early warning systems; and, as
    interventions of last resort, translocation or lethal removal \cite{muliya2026}. Each approach
    carries well-documented limitations that motivate the need for novel technological solutions.

    Physical barriers are expensive to install and maintain and are subject to habituation,
    particularly by intelligent megafauna such as elephants. In Kenya, the government has
    extended wildlife fencing by 200\,km as part of a recent operational response to HWC,
    while committing KES\,250\,million to the renovation of Meru Mulika Airstrip to boost
    regional tourism as a strategy for creating economic incentives for conservation
    \cite{kenya2025}. However, fencing cannot address the behaviour of habitual trespassers that
    have already learned to exploit human settlements, as these animals actively seek out
    gaps and weaknesses in barrier systems. Evidence from Laikipia County, where large-scale
    private ranches secured by electrified fencing reported fewer HWC incidents than the
    unfenced Kajiado landscape, supports the effectiveness of physical barriers where they can
    be fully maintained; however, this model is not scalable across most community and
    rangeland environments \cite{manoa2020}.

    Compensation schemes address economic grievances but do not prevent recurrence, and
    their effectiveness is severely undermined when payments are delayed or denied. In Kenya,
    the government's own acknowledgement that victims waited up to eight years for
    compensation between 2013 and 2022 illustrates the systemic dysfunction of previous
    arrangements \cite{kenya2025}. Research confirms that delayed or unpaid compensation
    significantly worsens community attitudes toward wildlife, generating hostility and
    increasing the likelihood of retaliatory killings \cite{basak2023,manoa2020}. In Kajiado and
    Laikipia Counties, 4,722 compensation claims were deferred nationally during the
    20142017 period due to missing documentation and insufficient funds, compounding
    community anger \cite{manoa2020}. Kenya's 2025 policy reforms - including the migration
    to a digital compensation platform, a 90-day reporting mandate, and a target of payment
    within four months - represent a significant institutional improvement \cite{kenya2025}, but
    do not address the upstream problem of preventing conflict incidents from occurring.

    Translocation of conflict-prone animals is extensively documented as problematic: elephants
    are highly intelligent, socially complex organisms whose translocation disrupts critical social
    networks, causes psychological distress, and frequently triggers homing behaviour, with
    displaced individuals navigating hazardous anthropogenic landscapes to return to their
    original ranges - effectively displacing rather than resolving conflict \cite{muliya2026}.
    Non-lethal deterrents such as beehive fences and chili fences have shown localised
    effectiveness \cite{mbiki2025}, but require sustained community maintenance and do not
    address the behaviour of confirmed habitual trespassers. Community education campaigns
    that highlight the ecological and economic benefits of wildlife, including the link between
    wildlife conservation and tourism revenue, have been effective in shifting perceptions in
    some contexts \cite{mbiki2025,basak2023}, but operate on timescales too slow to address
    immediate conflict events.

    Chemical immobilization remains the standard intervention for the safe capture of conflict
    animals, using dart projectors, blowpipe CO\textsubscript{2}-powered rifles, and projectile
    syringes loaded with veterinary sedatives such as etorphine hydrochloride or xylazine
    hydrochloride. The Kodagu programme employed trained veterinarians using etorphine
    hydrochloride (711\,mg depending on age and sex) facilitated by trained kumki elephants
    for safe positioning during sedation \cite{muliya2026}. While effective when conditions
    permit, this approach is limited by the logistical prerequisites of deploying trained personnel
    with specialised equipment in often remote terrain, close to a potentially aggressive animal.
    Response delays ranging from several hours to several years are common when managing
    multiple conflict individuals, and throughout these delays, casualties accumulate on both
    sides.

    \subsection{GPS Satellite Telemetry, Low-Cost Tracking, and Early Warning Systems}

    Satellite telemetry has substantially enhanced the capacity of wildlife managers to monitor
    conflict-prone animals in near real time and disseminate pre-emptive alerts to communities.
    Muliya et al.\ describe a landmark Early Warning System (EWS) implemented by the Wildlife
    Institute of India (WII) and the Karnataka Forest Department (KFD) in the coffee
    agroforestry landscapes of Kodagu district \cite{muliya2026}. Between 2019 and 2022,
    25 advanced satellite collars equipped with Iridium satellite transmitters were deployed
    on conflict-prone elephants, including six matriarchs representing herds of approximately
    90\,individuals and five lone tuskers responsible for known human fatalities. These
    programmable collars transmitted location data at frequencies of 5\,minutes to 24\,hours,
    enabling wildlife managers to track elephant movements in near real time.

    The resulting EWS disseminated telemetry data through multiple channels: rapid response
    teams for immediate conflict mitigation, local community alert networks, and research
    biologists for ecological study \cite{muliya2026}. Communities that previously advocated for
    elephant removal began actively supporting collaring programmes. Real-time tracking
    enabled coordinated Rapid Response Teams (RRTs) to drive elephants away from
    settlements, enforce temporary activity restrictions, and facilitate preventive evacuations.
    The programme accumulated approximately 104,764 elephant location points over
    4,516\,observation days, enabling predictive conflict risk mapping \cite{muliya2026}.
    However, the system retains a critical dependency on human intervention: detection of
    boundary crossing must be followed by mobilisation and deployment of field personnel before
    any physical intervention occurs. This step from detection to immobilization remains entirely
    dependent on human response time - the gap this research seeks to close.

    A parallel body of work has demonstrated that the cost barrier to GPS-based wildlife
    tracking can be substantially reduced without sacrificing positional accuracy. Im\-a\-mo\u{g}lu
    and Ta\c{s} developed an Arduino-based GPS tracking collar (FiT-SMART Collar\,2.0) at a
    total hardware cost of \euro{}92 - compared to approximately \euro{}400 for commercial
    equivalents such as the Garmin T5 \cite{imamoglu2022}. The device integrates a GPS module,
    GSM data transmission, MicroSD local storage, and a real-time clock module for
    configurable recording intervals, achieving positional accuracy within 23\,m of an
    industrial RTK-GPS benchmark during vehicle-based trials \cite{imamoglu2022}. Crucially,
    the GSM module enables real-time location data to be received as a Google Maps link via
    SMS from any location with cellular coverage, making the platform practically deployable
    in the field conditions typical of Kenyan rangeland counties. The authors identify
    battery life - reduced to approximately 48\,hours under continuous GSM operation -
    as the primary engineering limitation for long-term field deployment, and propose
    a remotely releasable collar design via servo motor as a future development
    \cite{imamoglu2022}. This low-cost platform is directly relevant to the proposed research:
    the GPS and GSM architecture of FiT-SMART Collar\,2.0 constitutes a validated, accessible
    foundation for the geofence trigger mechanism of the proposed immobilization device, and
    demonstrates that such technology can be made economically viable for wildlife agencies
    in lower-income settings.

    Related work in Kenya has shown that GSM-enabled radio collars can notify rangers in near
    real time when elephants cross into human-occupied areas, with approximately three-quarters
    of users taking pre-emptive actions upon receiving alerts (Graham et al., 2012, as cited in
    \cite{muliya2026}). In Sri Lanka, elephants have been fitted with aversive geofencing devices
    to manage their movements within virtual boundaries, representing the closest existing
    wildlife-domain analogue to the proposed device, though relying on aversive conditioning
    rather than pharmacological immobilization (de Mel et al., 2023, as cited in
    \cite{muliya2026}).

    \subsection{Virtual Fencing: Technology, Mechanisms, and Conservation Applications}

    Virtual fencing refers to GPS-enabled collar systems that enforce spatial boundaries for
    free-ranging animals using graduated sensory stimuli rather than physical barriers.
    Originally developed for livestock management, virtual fencing has evolved substantially
    and is now receiving growing recognition as a conservation tool across rangelands
    \cite{gilbert2026}. The technology enforces digital boundaries using GPS collars that
    deliver a warning sound when boundaries are approached, backed by a mild physical
    stimulus if the sound is ignored \cite{gilbert2026}. The standard design uses
    radio-frequency base stations or cellular networks to transmit collar locations to a digital
    platform via satellites, where animals' positions are precisely tracked and automatically
    compared to virtual boundaries. Warning sounds are delivered with directional cues so
    animals can ``hear'' which side the virtual fence is on, analogous to spatial orientation
    within herds \cite{gilbert2026}.

    A key feature relevant to the proposed research is the training period during which animals
    learn to respond to the warning sound alone, reducing the frequency of physical stimulus
    delivery. Research has found that stress responses of livestock to virtual fencing collars
    are low and comparable to those of physical fences, with precision training schedules
    ensuring animals rarely experience physical stimuli \cite{gilbert2026}. This demonstrated
    biocompatibility and behavioural efficacy of GPS collar-based management devices is
    directly informative for the proposed sedative delivery device.

    Gilbert identifies four primary current conservation functions of virtual fencing:
    substituting for physical fences to improve landscape connectivity; concentrating livestock
    to enhance habitat and reduce wildlife conflict; timing livestock rotations to improve habitat
    heterogeneity; and excluding livestock from sensitive areas \cite{gilbert2026}. Two further
    potential applications are identified: improving administration of conservation grazing
    programmes and rewarding ``conservation mode'' behaviour within software platforms. The
    technology has been piloted at multiple locations across the western United States, with
    projects testing livestock-carnivore conflict reduction, protection of riparian corridors,
    and invasive plant management.

    Particularly relevant for the proposed research, Gilbert notes that virtual fencing collars
    could enable ranchers to move cattle away from conflict risk areas when large carnivores
    are nearby, and that camera traps with edge computing are now capable of near real-time
    species detection that could trigger virtual fence boundary changes \cite{gilbert2026}. This
    convergence of geofencing with conflict risk detection is the direct conceptual antecedent
    of the proposed wildlife immobilization device, in which GPS geofence crossing by a conflict
    animal would trigger autonomous sedative delivery rather than a deterrent stimulus. Gilbert
    also documents limitations informative for device design: GPS accuracy can degrade in
    densely vegetated terrain, increasing animal welfare concerns; device retention across
    varying body conditions must be considered; and fail-safe mechanisms are essential to
    prevent inadvertent discharge \cite{gilbert2026}.

    \subsection{BioMEMS Devices: Materials, Architectures, and Implantable Drug Delivery}

    The field of Biomedical Microelectromechanical Systems (BioMEMS) encompasses
    microscale systems that integrate biosensors, actuators, microfluidics, wireless
    communication, and signal processing within compact, biocompatible platforms for use in
    medical applications \cite{abhinav2025}. The BioMEMS industry is projected to reach
    USD\,24.5\,billion by 2030, driven by rapid advances in miniaturised biomedical devices and
    digital health platforms \cite{abhinav2025}. BioMEMS devices broadly fall into two
    categories: wearable systems for non-invasive external monitoring, and implantable systems
    that work within the body for continuous internal sensing, targeted drug delivery, and
    neurostimulation \cite{abhinav2025}. For the proposed wildlife device, implantable BioMEMS
    architectures are of primary relevance.

    The selection of appropriate materials is fundamental to implantable BioMEMS performance
    and longevity. Abhinav et al.\ provide a comprehensive taxonomy of materials for BioMEMS
    applications, covering synthetic polymers, biodegradable and bioresorbable materials,
    natural polymers, and hybrid materials \cite{abhinav2025}. Polydimethylsiloxane (PDMS) is
    noted for its outstanding chemical stability, biocompatibility, and optical transparency,
    making it well suited for flexible substrates in implantable sensors. Polyimide offers
    exceptional mechanical strength (tensile strength $78.3 \pm 7.3$\,MPa) combined with
    high-temperature stability and excellent electrical insulation, making it ideal for neural
    electrodes and structural components in implantable devices \cite{abhinav2025}. Parylene\,C
    provides low water absorption (0.1\,\% after 24\,hours), high biocompatibility and
    hemocompatibility, and can be deposited as ultra-thin, defect-free protective coatings for
    implantable electronics \cite{abhinav2025}. Biodegradable polymers such as polylactic acid
    (PLA) and poly(lactic-co-glycolic acid) (PLGA) are of particular interest for drug delivery
    applications, as they degrade safely in vivo and can be formulated to control drug release
    kinetics over timescales ranging from days to months \cite{abhinav2025}. These material
    considerations directly inform the enclosure design, drug reservoir, and coating strategy
    for the proposed wildlife immobilization device.

    Implantable drug delivery systems represent one of the most clinically advanced
    applications of BioMEMS. The concept of implantable drug delivery was first introduced
    by Blackshear (1979), who explored subcutaneous drug-releasing implants and infusion
    pumps for continuous release, as documented by Abhinav et al.\ \cite{abhinav2025}. Since
    then, significant advances in microfabrication and MEMS technology have enabled the
    development of highly sophisticated implantable drug delivery systems (IDDSs). Abhinav
    et al.\ document a range of commercially deployed and research-stage systems, including
    the Medtronic Synchromed\,II programmable pump for chronic pain and spasticity
    management; MicroCHIPS Biotech MEMS-based microchip implants for wireless drug
    delivery; and research platforms combining nanochannel delivery membranes with
    Bluetooth Low Energy (BLE) control for chronic disease management \cite{abhinav2025}.
    Key challenges for implantable drug delivery identified by Abhinav et al.\ include energy
    supply (battery life and recharging), long-term biocompatibility against biofouling and
    inflammatory responses, regulatory complexity, data security, and device
    miniaturization - all of which must be addressed in the design of the proposed wildlife
    device \cite{abhinav2025}.

    The integration of BioMEMS with artificial intelligence (AI) and the Internet of Things
    (IoT) is further enhancing the field by enabling closed-loop, data-driven therapeutic
    delivery that responds to real-time physiological conditions \cite{abhinav2025}. AI algorithms
    can analyse physiological data to predict adverse events and personalise treatment
    regimens, while IoT connectivity enables remote monitoring and real-time alerts. For the
    proposed wildlife device, this paradigm translates directly: GPS location data from the
    implanted device can be relayed wirelessly to a ranger management platform, enabling
    real-time tracking of immobilized animals and integration with broader landscape-level
    HWC management systems.

    \subsection{Implantable Drug Delivery: Biocompatibility Evidence from Animal Models}

    The specific technical feasibility of an implantable, electrolytically actuated drug delivery
    device in a large animal model has been rigorously demonstrated by Guti\'errez-Hern\'andez
    et al.\ in a one-year feasibility study of the Replenish Posterior MicroPump \cite{gutierrez2014}.
    This novel bioelectronic MEMS device was developed for intravitreal drug delivery in
    chronic ocular disease and exemplifies the class of compact, battery-powered, programmable
    drug reservoirs that form the technical core of the proposed wildlife device. The MicroPump
    houses a rechargeable battery, a drug reservoir with a refill port, electronics, and an
    electrolysis chamber. Upon activation, electrolytic water splitting generates hydrogen and
    oxygen gas, and the resulting pressure deforms a membrane, forcing the drug from the
    reservoir through a cannula into the target site at a precisely preprogrammed microdose
    volume. The device is designed to last more than five years before requiring replacement
    \cite{gutierrez2014}.

    The feasibility study enrolled thirteen beagle dogs (11\,experimental, 2\,control), following
    animals for up to one year post-implantation. Ophthalmic examination and imaging revealed
    no clinically significant pathological changes, no significant inflammatory reaction, and no
    tissue ingrowth at the sclerotomy site \cite{gutierrez2014}. Histological evaluation showed
    only mild chronic inflammatory cell infiltration and a thin fibrous capsule around the device,
    consistent with the standard foreign body response to any implanted device and comparable
    to findings with established glaucoma drainage devices. No retinal toxic changes, no retinal
    pigment epithelium degeneration, and no evidence of fibrous ingrowth at the cannula site
    were observed. These findings confirm that an electrolytic MEMS micropump - with a
    biocompatible polymer enclosure, drug reservoir, pressure-driven delivery mechanism, and
    percutaneous cannula - can be maintained safely in a large animal for one year. While
    designed for a specialised ocular application, the core engineering principles -
    miniaturisation, programmable activation, electrolytic pumping, and biocompatible materials
    - are directly transferable to a subcutaneous sedative delivery device for wildlife
    management.

    \subsection{Fully Implantable Drug Delivery with Non-Invasive Refilling: The PILLSID System}

    A critical practical challenge for a long-term wildlife immobilization device is the need to
    replenish the drug reservoir without requiring surgical re-capture of the animal or
    explantation of the device. Iacovacci et al.\ address this challenge directly in their
    description of PILLSID (PILl-refiLled implanted System for Intraperitoneal Delivery), a
    fully implantable robotic device for programmable intraperitoneal drug delivery that is
    refilled through ingestible magnetic capsules \cite{iacovacci2021}. PILLSID was designed to
    restore blood glucose homeostasis in patients with type\,1 diabetes, but its engineering
    principles have broad applicability to any scenario requiring long-term implantable drug
    delivery with periodic reservoir replenishment.

    PILLSID is implanted in an extraperitoneal pouch fashioned between the costal margin and
    the iliac crest and delivers the target drug through an intraperitoneal catheter. The
    refilling procedure involves the oral ingestion of a drug-loaded magnetic capsule, which is
    passively transported along the gastrointestinal system until it reaches the implant and is
    magnetically docked by a magnetic switchable device (MSD). Once docked, a retractable
    needle punctures the capsule and transfers the drug to the implanted reservoir; the capsule
    is then released and naturally excreted \cite{iacovacci2021}. Once refilled, the device acts as
    a programmable microinfusion pump for precise drug release, operated via a Bluetooth
    communication protocol and a wirelessly rechargeable battery at the 13.56\,MHz medical
    band. The complete prototype measures $78 \times 63 \times 35$\,mm and weighs 165\,g,
    comparable in form factor to commercial implantable devices \cite{iacovacci2021}.

    In vivo validation in a diabetic swine model demonstrated that the refilling mechanism
    worked efficiently and that PILLSID could safely regulate blood glucose levels through
    intraperitoneal insulin delivery \cite{iacovacci2021}. The device achieved a microinfusion
    resolution of $3.18 \pm 0.30$\,$\mu$l regardless of time delay between consecutive injections
    across repeated emptying cycles, confirming dose precision. Wireless power transfer
    efficiency was demonstrated at $44 \pm 7$\,\% during in vivo validation in a curved
    configuration, confirming the feasibility of transcutaneous recharging \cite{iacovacci2021}.
    These results demonstrate that a fully implantable programmable drug pump can be
    maintained in a large animal, wirelessly recharged, and non-invasively refilled without
    surgical re-entry - all capabilities directly relevant to a wildlife device that must remain
    implanted across months or years of field deployment.

    For the proposed wildlife immobilization device, the PILLSID paradigm suggests a viable
    strategy for drug reservoir replenishment: a veterinarian could administer a drug-loaded
    capsule to a collared animal (either orally or via a specially designed delivery system) to
    refill the implanted sedative reservoir without requiring full anaesthetic recapture. This
    would substantially extend the operational life of the device and reduce the logistical
    burden of maintenance, particularly for wide-ranging conflict animals.

    \subsection{Autonomous Closed-Loop Implantable Drug Delivery: The Naloximeter Paradigm}

    The most directly analogous precedent for the proposed wildlife immobilization device is
    the Naloximeter, described by Ciatti et al.\ in \textit{Science Advances} \cite{ciatti2024}.
    The Naloximeter is a class of subcutaneously implanted devices developed to autonomously
    detect opioid overdose and deliver naloxone (NLX) without bystander intervention. Opioid
    overdose accounts for nearly 75,000 deaths per year in the United States, and its lethality
    depends on opioid-induced respiratory depression progressing to fatal hypoxia within
    minutes \cite{ciatti2024}. The Naloximeter circumvents the bystander requirement by
    detecting overdose physiologically and executing a closed-loop rescue cascade autonomously.

    The Naloximeter uses a dual-wavelength optical sensor (660\,nm and 880\,nm) to measure
    tissue oxygenation (StO\textsubscript{2}) as a proxy for respiratory status. A multivariable
    Overdose Detection Algorithm (ODA) analyses data - including rate of change of
    oxygenation, optical absorption, and motion index - to detect overdose within 90\,seconds
    of fentanyl administration and reject false positives during normal ambulation \cite{ciatti2024}.
    Upon detection, the rescue cascade activates the drug delivery mechanism and relays an
    emergency alert via Bluetooth Low Energy (BLE) and cellular communication to first
    responders, with a user abort option as a fail-safe against false positives.

    Two drug delivery embodiments were developed and tested. An intravenous version uses
    electrolytic pumping (functionally analogous to the MicroPump) to drive NLX through an
    indwelling catheter into the vasculature within 6\,minutes. An injector version uses a
    motor-driven plunger to deliver NLX subcutaneously within 25\,seconds, with dual-injection
    capability given that the potency of fentanyl analogues often requires multiple naloxone
    doses \cite{ciatti2024}. Pharmacokinetic data from porcine models showed NLX appearing in
    plasma within 1\,minute of pump activation. Survival studies in freely moving porcine models
    demonstrated rescue from otherwise fatal fentanyl overdose with recovery to baseline
    oxygenation within approximately 4.510\,minutes, compared to 23\,minutes for
    self-recovery ($P < 0.0001$) \cite{ciatti2024}. The device showed no significant biomarkers
    of cardiac damage post-activation, and catheters implanted for several weeks showed stable
    pharmacokinetics comparable to infusion pump controls. Battery life was approximately
    6.9\,days with wireless recharging demonstrated in porcine models.

    For the proposed wildlife immobilization device, the Naloximeter provides the complete
    template for closed-loop sensor-triggered autonomous drug delivery. The critical adaptation
    is the substitution of the physiological overdose sensor with a GPS geofence trigger:
    rather than monitoring tissue oxygenation for signs of respiratory depression, the device
    would monitor the animal's GPS coordinates relative to a pre-defined virtual boundary and
    initiate drug delivery upon crossing. This substitution is technically straightforward -
    compact GPS modules suitable for wildlife collar deployment are commercially available at
    low cost, with demonstrated positional accuracy of 23\,m \cite{imamoglu2022} - and is
    the conceptually novel contribution of the proposed research. Veterinary sedatives with
    well-characterised pharmacokinetic profiles (e.g., etorphine, medetomidine, or tiletamine-zolazepam combinations) would replace naloxone as the pharmacological
    payload, dosed by species and estimated body mass. The Naloximeter also establishes the
    viability of the emergency alert function, which in the wildlife context would notify ranger
    teams of an immobilized animal's location, enabling safe recovery and transport.

    \subsection{Synthesis: Converging Technologies and the Research Gap}

    A review of the literature across eight domains - HWC scale and
    community impacts in Kenya and globally
    \cite{muliya2026,manoa2020,mbiki2025,basak2023}, government policy and compensation
    frameworks \cite{kenya2025}, GPS-enabled virtual geofencing for animal management
    \cite{gilbert2026}, low-cost GPS tracking technology \cite{imamoglu2022}, BioMEMS
    materials and device architectures \cite{abhinav2025}, implantable MEMS micropump
    biocompatibility \cite{gutierrez2014}, non-invasive refilling of implantable drug delivery
    systems \cite{iacovacci2021}, and autonomous closed-loop drug delivery \cite{ciatti2024}
    - reveals a clear convergence of enabling technologies. When integrated, these advances
    could address the critical gap in current wildlife conflict management: the absence of an
    autonomous, immediate immobilization mechanism at the moment of boundary crossing.

    The empirical HWC record from Kenya establishes the urgency and species-specific
    parameters of the problem \cite{manoa2020}. Elephant-dominated conflict, peaking in dry-season
    months, produces the most severe human casualties and economic losses; yet these same
    animals - because of their intelligence and large home ranges - are the most resistant to
    conventional deterrence and the most logistically challenging to dart-tranquilize in a timely
    manner. Community perception data establish that the speed and reliability of any response
    mechanism is itself a determinant of conservation outcomes: communities that experience
    rapid, credible intervention are more likely to support wildlife conservation \cite{basak2023},
    and government investment in faster, more transparent HWC response systems reflects this
    political reality \cite{kenya2025}.

    The Kodagu EWS demonstrates that GPS satellite telemetry reliably tracks habitual
    trespassers in complex landscapes, that satellite collars can be safely applied to conflict
    megafauna, and that data-driven proactive management reduces HWC severity \cite{muliya2026}.
    The FiT-SMART Collar platform demonstrates that the GPS and GSM architecture required
    to implement geofence triggering can be assembled at a cost accessible to under-resourced
    wildlife agencies \cite{imamoglu2022}. Virtual fencing demonstrates that GPS collar devices
    can enforce spatial boundaries with graduated stimuli, that animals adapt to these devices
    safely, and that the same GPS architecture used for tracking can trigger interventions rather
    than merely alerts \cite{gilbert2026}. The BioMEMS review establishes the broad material,
    energy, and engineering foundations - biocompatible polymers, programmable
    microactuators, wireless power transfer, closed-loop architectures - that make
    sophisticated miniaturised implantable devices feasible for long-term deployment
    \cite{abhinav2025}. The MicroPump study establishes that an electrolytic MEMS drug pump
    can be safely maintained in a large animal for one year with minimal adverse effects
    \cite{gutierrez2014}. PILLSID demonstrates that a fully implantable programmable pump can
    be wirelessly recharged and non-invasively refilled without surgical re-entry - a critical
    operational requirement for a device deployed in free-ranging wildlife \cite{iacovacci2021}.
    The Naloximeter establishes conclusively that a closed-loop, sensor-triggered, implantable
    drug delivery system can autonomously rescue a freely moving large animal from a
    life-threatening event within minutes \cite{ciatti2024}.

    Yet no existing study or device combines GPS geofence boundary-detection with autonomous
    implantable sedative delivery for wildlife conflict management. Current GPS EWS programmes
    detect boundary crossing but cannot initiate immobilization autonomously. Virtual fencing
    deters animals but does not immobilize confirmed trespassers. Existing IDDSs are
    programmed for fixed schedules or physiological triggers, not geospatial triggers. The
    proposed research closes this gap by engineering a device at the intersection of these
    converging technologies: a subcutaneously implantable, GPS-triggered, electrolytically
    actuated sedative delivery system for habitual wildlife trespassers. The successful
    realisation of this device would represent a paradigm shift from reactive to proactive HWC
    management, eliminating the most dangerous phase of current practice : the delay between
    initial encroachment and ranger arrival : and thereby reducing human casualties, crop
    damage, animal deaths, and the erosion of conservation support in conflict-affected
    communities.

    % ══════════════════════════════════════════════════════════════════════════════


    % ══════════════════════════════════════════════════════════════════════════════
    % Replace the old thebibliography block with these two lines:
    \bibliographystyle{unsrtnat}
    \bibliography{references}

\end{document}